\chapter{Prompt and Skill Examples}

This appendix gives representative examples of the intended benchmark pressure. The exact prompt files are stored under \path{skill_benchmark/prompts/} and \path{skill_benchmark/prompts_public_gold/}.

\section{PDF Table Extraction}

\textbf{User request:} extract a table from a PDF while preserving row and column structure.

\textbf{Gold skill:} PDF layout table extractor.

\textbf{Plausible alternatives:} PDF question answerer, PDF OCR cleaner, document field extractor, PDF-to-DOCX converter.

\textbf{Procedural distinction:} the gold skill must inspect layout and reconstruct table structure. A PDF question-answering skill may ground answers in pages but not reconstruct tables. OCR cleaning handles scanned text uncertainty. Conversion changes format but is not optimized for table reconstruction.

\section{GitHub CI Root Cause}

\textbf{User request:} inspect a failed CI run and identify the likely root cause from logs.

\textbf{Gold skill:} CI log root-cause debugger.

\textbf{Plausible alternatives:} deployment build triager, PR review-comment resolver, repository code reviewer, release changelog generator.

\textbf{Procedural distinction:} the gold skill must parse CI logs and isolate the failing command, error, dependency, or environment condition. Review-comment resolution and code review are semantically adjacent GitHub workflows but require different inputs and outputs.

\section{Observability Trace Diagnosis}

\textbf{User request:} diagnose latency across a distributed trace path.

\textbf{Gold skill:} distributed trace investigator or implicit trace path diagnoser, depending on the prompt stratum.

\textbf{Plausible alternatives:} latency anomaly detector, SLO breach checker, service dependency mapper, metrics overview.

\textbf{Procedural distinction:} trace diagnosis follows spans and service boundaries. A metrics overview summarizes signals, and a dependency mapper describes architecture, but neither necessarily reconstructs a trace path.

\section{Public-Gold Example: Figma Implementation}

\textbf{User request:} implement an existing Figma node as production UI code with close visual parity.

\textbf{Gold skill:} public Figma implement-design skill.

\textbf{Plausible alternatives:} broad Figma parent skill, Figma generate-design skill, generic frontend implementation skill.

\textbf{Procedural distinction:} the gold skill targets an existing Figma node and produces code. A broad parent skill may route to several subskills, while generate-design writes or creates design content rather than implementing an existing node.

\section{Public-Gold Ambiguity Example}

Some public-gold prompts are not stable strict targets. For example, a public OCR skill may be valid for scanned PDF OCR, but a controlled OCR-cleaning skill may be better if the prompt asks for uncertain recognition regions and page anchors. These cases are either revised, excluded, or marked with acceptable alternatives before final reporting.
